\subsection{Applications: I. Extension of a continuous function to a space of measures}

Let T be a locally compact space, F a quasi-complete, Hausdorff locally convex space, and f a continuous mapping of T into F; if \(\mu\) is a positive measure on T, with compact support S, then \(\mathbf{f}(S)\) is compact; the closed convex envelope of \(\mathbf{f}(S)\) is then compact (TVS, III, §1, No.~6), therefore f is scalarly \(\mu\) -integrable and \(\int f d\mu \in F\) (No.~2, Cor. of Prop. 5). If now \(\lambda\) is any real measure with compact support, \(\lambda^{+}\) and \(\lambda^{-}\) are positive measures with compact support; setting \(\int f d\lambda = \int f d\lambda^{+} - \int f d\lambda^{-}\) , one verifies immediately (using the relation \((\lambda + \mu)^{+} + \lambda^{-} + \mu^{-} = \lambda^{+} + \mu^{+} + (\lambda + \mu)^{-}\) ) that \(\lambda \mapsto \int f d\lambda\) is a linear mapping of the space \(\mathcal{C}'(T)\) of measures on T with compact support, into the locally convex space F.

Let us now observe that the space \(\mathcal{C}^{\prime}(\mathrm{T})\) may be identified with the dual of the space \(\mathcal{C}(\mathrm{T})\) of continuous numerical functions on T (whence its notation), when \(\mathcal{C}(\mathrm{T})\) is equipped with the topology of compact convergence (which we shall always assume in this No.~and in the following one): for, it is known on the one hand (Ch. IV, §4, No.~8, Prop. 14) that the measures on T that can be extended to continuous linear forms on \(\mathcal{C}(\mathrm{T})\) are the measures with compact support, and conversely, the restriction to \(\mathcal{K}(T)\) of a continuous linear form on \(\mathcal{C}(T)\) is a measure (since the topology of \(\mathcal{K}(T)\) is finer than the one induced by the topology of \(\mathcal{C}(T)\) ).

PROPOSITION 14. --- Let T be a locally compact space, F a quasi-complete, Hausdorff locally convex space, and f a continuous mapping of T into F. If the space \(\mathcal{C}'(T)\) of measures on T with compact support is equipped with the topology of uniform convergence in the compact subsets of \(\mathcal{C}(T)\) , then the mapping \(\lambda \mapsto \int \mathbf{f} d\lambda\) is the unique continuous linear mapping \(\widetilde{\mathbf{f}}\) of \(\mathcal{C}'(T)\) into F such that \(\widetilde{\mathbf{f}}(\varepsilon_t) = \mathbf{f}(t)\) for every \(t \in T\) .

To establish the uniqueness of the extension, it suffices to see that the point measures \(\varepsilon_{t}\) form a total set in \(\mathcal{C}'(T)\) ; since the dual of \(\mathcal{C}'(T)\) is \(\mathcal{C}(T)\) (TVS, IV, §1, No.~1, Th. 1), it suffices to observe that every function \(g \in \mathcal{C}(T)\) that is orthogonal to all of the measures \(\varepsilon_{t}\) is equal to 0 by definition (TVS, IV, §1, No.~2, Prop. 2).

Let us now show that \(\lambda \mapsto \int \mathbf{f}d\lambda\) is continuous. Let \(V\) be a closed, balanced convex neighborhood of 0 in \(F\) ; it suffices to prove that there exists a relatively compact subset \(L\) of \(\mathcal{C}(T)\) such that the relations \(\lambda \in L^{\circ}\) and \(\mathbf{z}' \in V^{\circ}\) imply \(\left|\left\langle \int \mathbf{f}d\lambda, \mathbf{z}' \right\rangle \right| \leqslant 1\) , or again that \(\left|\int \langle \mathbf{f}, \mathbf{z}' \rangle d\lambda \right| \leqslant 1\) . To this end, we are going to show that as \(\mathbf{z}'\) runs over \(V^{\circ}\) , the set \(L\) of numerical functions \(\langle \mathbf{f}, \mathbf{z}' \rangle\) is relatively compact in \(\mathcal{C}(T)\) . Since \(V^{\circ}\) is bounded for \(\sigma(F', F)\) , the supremum of the numbers \(|\langle \mathbf{f}(t), \mathbf{z}' \rangle|\) , for \(t \in T\) fixed and \(\mathbf{z}'\) running over \(V^{\circ}\) , is finite; by virtue of Ascoli's theorem (GT, X, §2, No.~5, Cor. 2 of Th. 2), it therefore suffices to show that the set of \(\langle \mathbf{f}, \mathbf{z}' \rangle\) ( \(\mathbf{z}' \in V^{\circ}\) ) is equicontinuous. But, for every \(t_0 \in T\) and every \(\delta > 0\) , there exists by hypothesis a neighborhood \(W\) of \(t_0\) in \(T\) such that \(\mathbf{f}(t) - \mathbf{f}(t_0) \in \delta V\) for all \(t \in W\) ; it follows that \(|\langle \mathbf{f}(t), \mathbf{z}' \rangle - \langle \mathbf{f}(t_0), \mathbf{z}' \rangle| \leqslant \delta\) for all \(t \in W\) and all \(\mathbf{z}' \in V^{\circ}\) , which completes the proof.

Remarks. --- 1) The mapping \(t \mapsto \varepsilon_t\) is a homeomorphism of T into the space \(\mathcal{C}'(\mathrm{T})\) ; for, if L is a compact subset of \(\mathcal{C}(\mathrm{T})\) , and \(t_0 \in \mathrm{T}\) , there exists (GT, X, §2, No.~5, Cor. 3 of Th. 2) a neighborhood W of \(t_0\) such that \(|g(t) - g(t_0)| \leqslant 1\) for every \(t \in \mathrm{W}\) and every function \(g \in \mathrm{L}\) , therefore \(\varepsilon_t - \varepsilon_{t_0} \in \mathrm{L}^\circ\) for \(t \in \mathrm{W}\) , which proves the continuity of \(t \mapsto \varepsilon_t\) (cf.~Ch. IV, §4, No.~8, Prop. 15); one knows in addition that the inverse mapping is already continuous for the vague topology (Ch. III, §1, No.~9, Prop. 13), hence a fortiori for the topology of uniform convergence in the compact subset of \(\mathcal{C}(\mathrm{T})\) . If one then identifies T with its image in \(\mathcal{C}'(\mathrm{T})\) via \(t \mapsto \varepsilon_t\) , one can say that \(\lambda \mapsto \int f d\lambda\) is the unique continuous extension of f to a linear mapping.

\begin{enumerate}
\def\labelenumi{\arabic{enumi})}
\setcounter{enumi}{1}
\tightlist
\item
  Note that in the proof of the continuity of \(\lambda\mapsto\int f d\lambda\) , we have not used the fact that F is quasi-complete. The conclusion of Prop. 14 is therefore still valid without this hypothesis, provided one knows in addition that \(\int f d\mu\in F\) for every positive measure \(\mu\) with compact support.
\end{enumerate}

Suppose now that \(\mathbf{f}(\mathrm{T})\) is a bounded subset of \(\mathrm{F}\) . Then, for every bounded positive measure \(\mu\) on \(\mathrm{T}\) , \(\mathbf{f}\) is scalarly \(\mu\) -integrable and \(\int \mathbf{f}d\mu \in \mathrm{F}\) (No.~2, Prop. 8). If \(\lambda\) is any bounded real measure on T, then \(\lambda^{+}\) and \(\lambda^{-}\) are bounded, and one sees immediately that \(\lambda \mapsto \int \mathbf{f} d\lambda\) , defined as above, is a linear mapping of the space \(\mathcal{M}^1(\mathrm{T})\) of bounded measures on T, into the locally convex space F, that obviously extends the mapping \(\lambda \mapsto \int \mathbf{f} d\lambda\) of \(\mathcal{C}'(\mathrm{T})\) into F.

PROPOSITION 15. --- Let T be a locally compact space, F a Hausdorff locally convex space that is quasi-complete, and f a continuous mapping of T into F such that \(\mathbf{f}(\mathrm{T})\) is bounded. If \(\mathcal{M}^{1}(\mathrm{T})\) is equipped with its Banach space topology, then the linear mapping \(\lambda \mapsto \int f d\lambda\) of \(\mathcal{M}^{1}(\mathrm{T})\) into F is continuous.

For every closed, balanced, convex neighborhood V of 0 in F, there exists a \(\rho > 0\) such that \(\mathbf{f}(\mathrm{T}) \subset \rho\mathrm{V}\) ; the closed, balanced, convex envelope B of \(\mathbf{f}(\mathrm{T})\) is therefore contained in \(\rho V\) , and it is complete by hypothesis. If then \(\|\lambda\| \leqslant 1/\rho\) , it follows from No.~2, Prop. 8, and the relation \(\|\lambda\| = \lambda^{+}(\mathrm{T}) + \lambda^{-}(\mathrm{T})\) , that \(\int f d\lambda \in B/\rho \subset V\) .

